% =====================================================================
\section{Conclusion}
\label{sec:conclusion}
% =====================================================================

This document has translated the LipScribe proposal into an implementation plan.
The system is built from separable modules: a visual stream that uses MediaPipe
Face Mesh to extract stabilised $96\times96$ mouth crops and encodes them with a 3D
convolutional front-end and a BiLSTM; an audio stream that converts 16\,kHz audio
into log-Mel features and encodes them with strided convolutions and a BiLSTM at
the same 25\,fps rate; a cross-modal attention and gating module that combines
the two; and a CTC output layer decoded by prefix beam search. Modality dropout
and multi-condition training with MUSAN noise are the mechanisms through which the
fusion is expected to learn to rely on the lips when the audio degrades.

The architecture, the development stages and the evaluation plan are tied
together by design. Each milestone produces a component with its own tests; the
baselines reuse the same components, so that every comparison isolates one
change; and the per-SNR evaluation on held-out speakers measures exactly the
property the project claims to improve. The same preprocessing code serves
training and the captioning demonstration, so the demonstrated behaviour reflects
the evaluated model.

The plan also states what remains uncertain. The clean-condition target of WER
below 10\% and the expected advantage at 0 and $-5$\,dB are hypotheses to be
tested, not results. Emotion detection and language identification require data
that GRID does not provide and are kept as optional extensions. By completing the
core recogniser first, reporting outcomes faithfully---including any in which
fusion does not help---and documenting every experiment, the team aims to deliver
a working, verifiable audio-visual captioning system within the six-week course.
